\section{Related Work}
\label{sec:related}

\paragraph{Language models as email and SMS detectors.}
Using pre-trained language models without fine-tuning (zero-shot classification)
is now a standard baseline for security triage. Rojas-Galeano tested models on
SpamAssassin without extra training. They reported an F1 score of 90\% for
Flan-T5 on short emails and 95\% for GPT-4 on summarized
emails~\cite{rojas2024zeroshot}. Eilertsen et al.\ tested zero-shot and few-shot
prompts for phishing detection. They found mixed accuracy across models and
questioned if decisions rely on model knowledge or prompt
examples~\cite{eilertsen2025intent}. Koide et al.\ designed a detector using
email headers and bodies. They reported 99.70\% accuracy with
GPT-4~\cite{koide2024chatspamdetector}. Wang et al.\ added external web context
to SMS triage inside an agent pipeline, achieving 98.8\% overall
accuracy~\cite{wang2025smishx}. Thapa et al.\ studied small language models and
observed that rewritten emails lower detection
performance~\cite{thapa2025quantized}. Our study focuses on measurement
discipline rather than model architecture. We keep the detector and prompt
constant, applying careful accounting rules to measure how text edits change the
verdict.

\paragraph{Attacking security text filters.}
Shafee et al.\ generated plausible text to attack threat intelligence
pipelines. They found that evasion misled classifiers, producing a 75\%
false-positive rate when ChatGPT-4o acted as the
detector~\cite{shafee2026falsealarms}. Kulkarni et al.\ generated adversarial
phishing pages. They observed that small detectors failed, while vision-language
models proved more resilient~\cite{kulkarni2024mltollm}. In our setting, we
study raw message text rather than web pages. Our adversary makes small edits
rather than generating new text from scratch. This matches the real constraints
of an attacker who wants users to read the message.

\paragraph{Text perturbation techniques and budgets.}
Character-level attacks change individual letters. DeepWordBug introduced typos
and letter swaps~\cite{gao2018deepwordbug}, while HotFlip formalized gradient
character flips~\cite{ebrahimi2018hotflip}. Word-level attacks replace whole
words while preserving meaning. TextFooler searches synonyms
greedily~\cite{jin2020textfooler}, and PWWS ranks words by probability changes
(saliency)~\cite{ren2019pwws}. The TextAttack library organizes these methods
into standard goals, constraints, and transformations~\cite{morris2020textattack}.
In this paper, we evaluate two main search strengths. The first is a random
draw over a small synonym table. This represents a simple adversary and serves
as a lower bound on evasion rates. The second is a guided attack that ranks word
swaps by their measured effect on the model (\S\ref{sec:saliency}). In
\S\ref{sec:ablation}, we also test TextFooler's leave-one-out word ranking as a
reference arm. We use PWWS and TextFooler to define search logic rather than to
compare numerical results directly.

\paragraph{Cross-model transfer and evaluation methodology.}
Yuan et al.\ showed that adversarial text often transfers across different
neural models, and they built ensembles to target multiple models
together~\cite{yuan2021transfer}. He et al.\ extracted model parameters to build
transferable attacks against a target system~\cite{he2021extraction}. These
studies motivate our cross-model transfer matrix in \S\ref{sec:transfer}. In our
tests, black-box perturbations found on one model rarely transfer to another.
Our adversary uses a bounded local search rather than learning global transfer
directions. On statistical methodology, Carlini et al.\ demonstrated that many
adversarial robustness evaluations contain flaws and proposed standard practices
for threat models~\cite{carlini2019evaluating}. Zimmermann et al.\ used active
tests to reveal weaknesses in 11 of 13 published defenses~\cite{zimmermann2022confidence}.
Recent research also highlights the importance of exact confidence intervals for
robustness certificates~\cite{voracek2024statistical}. We follow these
rigorous practices. Every evasion rate in this paper reports its numerator, its
denominator, and an exact Wilson interval~\cite{nistwilson}.
